\documentclass{stex}
\libusepackage{naproche}
\libusepackage{copyright}
\libinput{preamble}

\title{``Little Gauß''' Theorem in \Naproche}
\author{Christian Schöner, Marcel Schütz}
\date{2025}

\begin{document}
\begin{smodule}{little-Gauss.ftl}
\maketitle

\importmodule[libraries/arithmetics]{multiplication.ftl}

\symdef{Gauss}{\mathcal{G}}
\symdecl*{Little Gauss}

\noindent This is a formalization of ``Little Gauß''' Theorem, i.e. of
the assertion that
\[\sum_{k = \Zero}^n k \Eq \frac{k(k \Plus \One)}\Two\]
for all $n \Elem \Naturals$.
In \Naproche we can define the sum $\sum_{n = \Zero}^k n$ by the function
$\Gauss : \Naturals \to \Naturals$ with $\Gauss(\Zero) \Eq \Zero$ and
$\Gauss(n \Plus \One) \Eq \Gauss(n) \Plus (n \Plus \One)$:

\begin{forthel}
  Let $n$ denote a natural number.

  \begin{signature*}[for=Gauss]
    $\Gauss(n)$ is a natural number.
  \end{signature*}

  \begin{axiom*}
    $\Gauss(\Zero) \Eq \Zero$.
  \end{axiom*}

  \begin{axiom*}
    $\Gauss(n \Plus \One) \Eq \Gauss(n) \Plus (n \Plus \One)$.
  \end{axiom*}
\end{forthel}

\noindent Then the theorem can be stated as follows.

\begin{forthel}
  \begin{theorem*}[title=Little Gauß,name=Little Gauss]
    For all natural numbers $n$ we have
    \[\Two \Mul \Gauss(n) \Eq n \Mul (n \Plus \One).\]
  \end{theorem*}
  \begin{proof}
    Define $\Phi \DefEq \SClass{n}{\Naturals}{\Two \Mul \Gauss(n) \Eq n \Mul (n \Plus \One)}$.
    
    (1) $\Zero \Elem \Phi$.

    (2) For all $n \Elem \Phi$ we have $n \Plus \One \Elem \Phi$.
    \begin{proof}
      Let $n \Elem \Phi$.
      Then $\Two \Mul \Gauss(n) \Eq n \Mul (n \Plus \One)$.
      Hence
      \begin{align*}
              \Two \Mul \Gauss(n \Plus \One)
        & \Eq \Two \Mul (\Gauss(n) \Plus (n \Plus \One)) \\
        & \Eq (\Two \Mul \Gauss(n)) \Plus (\Two \Mul (n \Plus \One)) \\
        & \Eq (n \Mul (n \Plus \One)) \Plus (\Two \Mul (n \Plus \One)) \\
        & \Eq (n \Plus \Two) \Mul (n \Plus \One) \\
        & \Eq (n \Plus \One) \Mul (n \Plus \Two) \\
        & \Eq (n \Plus \One) \Mul ((n \Plus \One) \Plus \One)
      \end{align*}.
    \end{proof}
  \end{proof}
\end{forthel}

\printbibliography
\printlicense[CcByNcSa]{Christian Schöner (2024), Marcel Schütz (2024--2025)}
\end{smodule}
\end{document}
